% ==============================================================================
% 07_criteria4_assessment.tex — Criterion 4: Student Assessment
% Score: 2
% ==============================================================================

\criterion{4}{การประเมินผู้เรียน}{Student Assessment}

% ------------------------------------------------------------------------------
\criterionsub{4.1 The student assessment is shown to use a variety of assessment methods that are constructively aligned with the intended learning outcomes.}

หลักสูตรกำหนดวิธีการประเมินที่หลากหลายสำหรับทุกรายวิชา โดยออกแบบให้สอดคล้องกับ CLOs ตามหลัก Constructive Alignment กล่าวคือ วิธีประเมินแต่ละชิ้นมีเป้าหมายในการวัด CLO ที่กำหนดไว้อย่างชัดเจน ดังแสดงในตาราง~\ref{tbl:assessment_methods}

\begin{table}[H]
  \centering
  \caption{วิธีประเมินและ CLO ที่วัดในวิชาบังคับหลัก}
  \label{tbl:assessment_methods}
  \begingroup
  \footnotesize
  \setlength{\tabcolsep}{3pt}
  \renewcommand{\arraystretch}{1.2}
  \begin{tabularx}{\linewidth}{|>{\raggedright\arraybackslash}X|>{\raggedright\arraybackslash}X|>{\raggedright\arraybackslash}X|}
    \hline
    \textbf{รายวิชา} & \textbf{วิธีประเมิน} & \textbf{CLO ที่วัด} \\
    \hline
    4095101 AI and ML & สอบข้อเขียน, โครงงาน DL, รายงาน, สาธิตระบบ, บทความสะท้อนคิด & CLO1, CLO2, CLO3, CLO4, CLO5 \\
    \hline
    4095102 Advanced Programming for ML & แบบทดสอบ, สอบปฏิบัติ Git/GitHub, รายงานวิเคราะห์ข้อมูล, โครงงาน deploy & CLO1--CLO5 \\
    \hline
    7015101 Emerging Tech & รายงาน, นำเสนอ IEL project, การมีส่วนร่วม workshop & CLO1--CLO5 \\
    \hline
    7015906 Research Methodology & สอบข้อเขียน, โครงร่างวิจัย, รายงานทบทวนวรรณกรรม, วิเคราะห์สถิติ & CLO1--CLO5 \\
    \hline
    7015907 Seminar & การนำเสนอบทความ, รายงาน, การมีส่วนร่วมอภิปราย & CLO1--CLO5 \\
    \hline
    7015901--7015903 วิทยานิพนธ์ 1--3 & ข้อเสนอโครงร่าง, รายงานความก้าวหน้า, สอบปากเปล่า, วิทยานิพนธ์ฉบับสมบูรณ์ & CLO1--CLO4 (แต่ละขั้น) \\
    \hline
  \end{tabularx}
  \endgroup
\end{table}

โครงสร้างคะแนนทุกวิชาบังคับกำหนดเป็นสัดส่วน 30\%/40\%/30\% (ก่อนกลางภาค/กลางภาค/ปลายภาค) เพื่อให้การประเมินกระจายตลอดภาคการศึกษา ไม่กระจุกอยู่ที่การสอบปลายภาคเพียงอย่างเดียว ซึ่งสะท้อนหลักการ Formative และ Summative Assessment ที่สมดุล

% ------------------------------------------------------------------------------
\criterionsub{4.2 The programme to show that its student assessment policies and regulations are explicit, fair, and consistently applied.}

มหาวิทยาลัยราชภัฏอุตรดิตถ์มีนโยบายการประเมินผู้เรียนและระบบการอุทธรณ์ที่กำหนดไว้เป็นลายลักษณ์อักษรและเผยแพร่อย่างเป็นทางการ

\textbf{นโยบายการประเมิน}: กำหนดเกณฑ์การประเมิน รูปแบบการให้คะแนน และกระบวนการตัดสินผลการเรียนไว้อย่างชัดเจนใน มคอ.3 ทุกรายวิชา นักศึกษาได้รับทราบนโยบายตั้งแต่ต้นภาคการศึกษา ผ่านการแจ้งใน Course Syllabus และการปฐมนิเทศรายวิชา

\textbf{ระบบการอุทธรณ์}: มหาวิทยาลัยออกประกาศระบบและแนวทางการอุทธรณ์ผลการประเมิน พ.ศ.~2567 (\hyperlink{ref:AUNQA-4-2-1}{AUNQA-4-2-1}) ซึ่งกำหนดขั้นตอนการอุทธรณ์ที่ชัดเจน ผู้รับผิดชอบในแต่ละขั้นตอน และระยะเวลาในการดำเนินการ โดยเผยแพร่ผ่านเว็บไซต์กองบริการการศึกษา มหาวิทยาลัยราชภัฏอุตรดิตถ์ นักศึกษาทุกคนสามารถเข้าถึงข้อมูลดังกล่าวได้ตลอดเวลา

ระบบนี้ใช้อย่างคงเส้นคงวาสำหรับนักศึกษาทุกหลักสูตรในมหาวิทยาลัย รวมถึงหลักสูตร วศ.ม. CPE\&AI

% ------------------------------------------------------------------------------
\criterionsub{4.3 The programme to show that the student progression and degree completion standards are explicit, communicated, and consistently applied.}

มาตรฐานและเกณฑ์ความก้าวหน้าทางการศึกษาสำหรับระดับบัณฑิตศึกษาของมหาวิทยาลัยราชภัฏอุตรดิตถ์กำหนดไว้ใน ข้อบังคับมหาวิทยาลัยราชภัฏอุตรดิตถ์ ว่าด้วยการศึกษาระดับบัณฑิตศึกษา พ.ศ.~2566 (\hyperlink{ref:AUNQA-4-3-1}{AUNQA-4-3-1}) และฉบับที่~2 (\hyperlink{ref:AUNQA-4-3-2}{AUNQA-4-3-2}) ครอบคลุมสาระสำคัญดังนี้

\begin{enumerate}
  \item \textbf{คุณสมบัติผู้เข้าศึกษา}: กำหนดวุฒิการศึกษา คะแนนเฉลี่ย และคุณสมบัติเฉพาะสำหรับ วศ.ม. CPE\&AI
  \item \textbf{เกณฑ์ความก้าวหน้าวิทยานิพนธ์}: นักศึกษาแผน~1 ต้องผ่านการเสนอหัวข้อ (วิทยานิพนธ์~1) ก่อนดำเนินการเก็บข้อมูล (วิทยานิพนธ์~2) และนำเสนอผลการวิจัย (วิทยานิพนธ์~3) ตามลำดับ
  \item \textbf{การพ้นสภาพนักศึกษา}: กำหนดเงื่อนไขการพ้นสภาพอย่างชัดเจน เช่น เรียนเกินระยะเวลาที่กำหนด ผลการเรียนต่ำกว่าเกณฑ์
  \item \textbf{เงื่อนไขสำเร็จการศึกษา}: ต้องสะสมหน่วยกิตครบ~36 หน่วยกิต (ขั้นต่ำ) ผ่านการสอบวิทยานิพนธ์/สารนิพนธ์ และได้รับการอนุมัติจากคณะกรรมการบัณฑิตศึกษา
\end{enumerate}

ข้อบังคับดังกล่าวเผยแพร่ผ่านเว็บไซต์มหาวิทยาลัยและสื่อสารแก่นักศึกษาในกิจกรรมปฐมนิเทศ นักศึกษาทุกคนสามารถตรวจสอบสถานะความก้าวหน้าของตนเองได้ผ่านระบบสารสนเทศของมหาวิทยาลัย

% ------------------------------------------------------------------------------
\criterionsub{4.4 The programme to show that its student assessment uses rubrics, marking schemes, and regulations that ensure validity, reliability, and fairness.}

หลักสูตรกำหนดเครื่องมือและกระบวนการประเมินที่รับประกันความตรง ความเที่ยง และความยุติธรรมในทุกรายวิชา

\textbf{Rubrics สี่ระดับคุณภาพ}: มคอ.3 วิชาบังคับทุกวิชากำหนด Analytic Scoring Rubric สำหรับ CLO แต่ละข้อ โดยระบุ~4 ระดับคุณภาพ (ดีเยี่ยม/ดี/ผ่าน/ต้องปรับปรุง) พร้อมคำพรรณนาที่ชัดเจนในแต่ละระดับ ซึ่งทำให้การให้คะแนนมีเกณฑ์ที่โปร่งใสและตรวจสอบได้

\textbf{Validity (ความตรง)}: วิธีประเมินออกแบบให้ตรงกับ CLO ที่ต้องการวัด เช่น CLO ที่ต้องการให้นักศึกษา ``พัฒนาระบบ AI'' วัดด้วยการส่งโครงงานและสาธิตระบบ (ไม่ใช่แค่การสอบข้อเขียน) CLO ``อธิบายหลักการ'' วัดด้วยข้อสอบข้อเขียนหรือการนำเสนอ

\textbf{Reliability (ความเที่ยง)}: หลักสูตรใช้แนวปฏิบัติและการกำกับติดตามการส่งผลการประเมิน พ.ศ.~2567 (\hyperlink{ref:AUNQA-4-7-1}{AUNQA-4-7-1}) เพื่อให้มั่นใจว่าการประเมินดำเนินการตรงตามกำหนดเวลา และใช้เกณฑ์ที่กำหนดไว้อย่างสม่ำเสมอ

\textbf{Fairness (ความยุติธรรม)}: หลักสูตรใช้การตัดเกรดแบบอิงเกณฑ์ (Criterion-Referenced) ทุกรายวิชา นักศึกษาทุกคนได้รับการประเมินด้วยเกณฑ์เดียวกัน และมีระบบการอุทธรณ์ผลประเมินรองรับกรณีที่นักศึกษาเห็นว่าไม่ได้รับความยุติธรรม

\textbf{Timelines}: กำหนดวันส่งงาน วันสอบ และกำหนดการคืนผลการประเมินไว้ใน มคอ.3 และแจ้งนักศึกษาตั้งแต่ต้นภาคการศึกษา

% ------------------------------------------------------------------------------
\criterionsub{4.5 The programme to show that student assessment is used to measure both the achievement of the intended course learning outcomes and the contribution of each course to the programme learning outcomes.}

หลักสูตรกำหนดระบบวัดการบรรลุ CLOs และ PLOs ดังนี้

\textbf{การวัดบรรลุ CLOs รายวิชา}: อาจารย์ผู้สอนทุกวิชาวิเคราะห์ผลการประเมินของนักศึกษาเทียบกับ Rubric ที่กำหนดไว้ในแต่ละ CLO หาก CLO ใดมีนักศึกษาไม่บรรลุเกณฑ์เกินร้อยละ~30 อาจารย์จะระบุในรายงาน มคอ.5 พร้อมแนวทางปรับปรุง

\textbf{Achievement Criteria สำหรับ CLOs}: กำหนดว่านักศึกษา ``บรรลุ'' CLO เมื่อได้คะแนนในระดับ ``ผ่าน'' ขึ้นไป (ระดับ~3 จาก~4) ในการประเมิน CLO นั้นๆ (เป้าหมาย: นักศึกษาไม่น้อยกว่า~80\% บรรลุ CLO แต่ละข้อ)

\textbf{ผลการวัดบรรลุ CLOs จริง (ปีการศึกษา~2568 — มคอ.5 ลงนาม 30 เมษายน~2569)}: หลักสูตรดำเนินการวัดและรายงานผล CLO Achievement ผ่านรายงาน มคอ.5 ครบทุกวิชาบังคับหลัก~4 วิชาแล้ว ดังแสดงในตาราง~\ref{tbl:clo_achievement_mko5}

\begin{table}[H]
  \centering
  \caption{ผลการสัมฤทธิ์ผลการเรียนรู้รายวิชา (CLO Achievement) จากรายงาน มคอ.5 ปีการศึกษา~2568}
  \label{tbl:clo_achievement_mko5}
  \begingroup
  \footnotesize
  \setlength{\tabcolsep}{3pt}
  \renewcommand{\arraystretch}{1.2}
  \begin{tabularx}{\linewidth}{|>{\raggedright\arraybackslash}X|c|c|c|c|c|c|c|}
    \hline
    \textbf{รายวิชา} & \textbf{n} & \textbf{CLO1} & \textbf{CLO2} & \textbf{CLO3} & \textbf{CLO4} & \textbf{CLO5} & \textbf{สัมฤทธิ์} \\
    \hline
    4095101 ปัญญาประดิษฐ์และการเรียนรู้ของเครื่อง (ภาค 1/2568) & 4 & 100\% & 100\% & 100\% & 100\% & 100\% & \textbf{100\%} \\
    \hline
    4095102 การเขียนโปรแกรมขั้นสูงสำหรับ ML (ภาค 1/2568) & 4 & 100\% & 100\% & 100\% & 100\% & 100\% & \textbf{100\%} \\
    \hline
    7015101 เทคโนโลยีอุบัติใหม่ฯ (ภาค 2/2568) & 4 & 100\% & 100\% & 100\% & 100\% & 100\% & \textbf{100\%} \\
    \hline
    7015906 ระเบียบวิธีวิจัยทางวิทยาศาสตร์และวิศวกรรมศาสตร์ (ภาค 2/2568) & 4 & 100\% & 100\% & 100\% & 100\% & 100\% & \textbf{100\%} \\
    \hline
    \multicolumn{2}{|l|}{\textbf{ภาพรวม (ทุกวิชา ทุก CLO)}} & \multicolumn{5}{c|}{\textbf{100\% ทุก CLO ทุกวิชา}} & \textbf{100\%} \\
    \hline
  \end{tabularx}
  \endgroup
  \begin{flushleft}
    \footnotesize เป้าหมาย: นักศึกษาไม่น้อยกว่า 80\% บรรลุ CLO แต่ละข้อ | ผลจริง: เกินเป้าหมาย 20 percentage points ทุกวิชา ทุก CLO \\
    อ้างอิง: มคอ.5 AUN3\_4095101, AUN3\_4095102, AUN3\_7015101, AUN3\_7015906 (ลงนาม 30 เมษายน 2569) — \hyperlink{ref:AUNQA-4-5}{AUNQA-4-5}
  \end{flushleft}
\end{table}

ผลการสัมฤทธิ์ผล~100\% สะท้อนให้เห็นถึงความสอดคล้องระหว่างการออกแบบการสอน (C3) การประเมิน (C4) และความสามารถของนักศึกษา ซึ่งในส่วน มคอ.5 ยังรายงาน ``ปัญหาและแผนปรับปรุง'' ไว้ด้วย เช่น วิชา 7015101 รายงานว่านักศึกษา~1 คนมีปัญหาด้านภาษาอังกฤษในการอ่านบทความ และแนะนำให้บูรณาการทักษะภาษาอังกฤษเชิงวิชาการเข้าไว้ในกิจกรรมสัมมนา ซึ่งแสดงให้เห็นว่า Assessment System สามารถตรวจพบจุดที่ต้องปรับปรุงได้จริง

\textbf{การวัดบรรลุ PLOs ระดับหลักสูตร}: หลักสูตรกำหนด PLO Achievement Criteria ระดับหลักสูตรดังตาราง~\ref{tbl:plo_achievement_criteria}

\begin{table}[H]
  \centering
  \caption{เกณฑ์การบรรลุ PLO ระดับหลักสูตร (PLO Achievement Criteria)}
  \label{tbl:plo_achievement_criteria}
  \begingroup
  \footnotesize
  \setlength{\tabcolsep}{3pt}
  \renewcommand{\arraystretch}{1.2}
  \begin{tabularx}{\linewidth}{|c|>{\raggedright\arraybackslash}X|>{\raggedright\arraybackslash}X|}
    \hline
    \textbf{PLO} & \textbf{เกณฑ์การบรรลุ (Achievement Criterion)} & \textbf{หลักฐาน} \\
    \hline
    PLO1 & นักศึกษาไม่น้อยกว่า~80\% บรรลุ CLO1/CLO2 ของรายวิชา AI/ML ในระดับ ``ผ่าน'' ขึ้นไป (ระดับ~3/4) และผ่านการสอบวิทยานิพนธ์/สารนิพนธ์ & รายงาน มคอ.5 + ผลการสอบป้องกัน \\
    \hline
    PLO2 & นักศึกษาไม่น้อยกว่า~80\% บรรลุ CLO ที่ map PLO2 ในระดับ ``ผ่าน'' ขึ้นไป; Community project ผ่านเกณฑ์ & รายงาน มคอ.5 + รายงาน IEL project \\
    \hline
    PLO3 & นักศึกษาทุกคนผ่านการสอบป้องกันวิทยานิพนธ์/สารนิพนธ์โดยคณะกรรมการ (คะแนนรวมระดับ ``ผ่าน'' ตามเกณฑ์ข้อบังคับ มรอ. 2566) & ผลการสอบปากเปล่า \\
    \hline
    PLO4 & นักศึกษาได้คะแนน CLO4 วิชา 7015906 ระดับ ``ผ่าน'' ขึ้นไป (Affective~5) และผ่านแบบประเมินจริยธรรมวิจัยก่อนสำเร็จการศึกษา & รายงาน มคอ.5 + แบบประเมิน \\
    \hline
    PLO5 & นักศึกษาได้คะแนน CLO5 วิชา 7015906 ระดับ ``ผ่าน'' ขึ้นไป และผ่าน Self-assessment LLL ก่อนสำเร็จการศึกษา & รายงาน มคอ.5 + แบบ Self-assessment \\
    \hline
  \end{tabularx}
  \endgroup
\end{table}

ในฐานะหลักสูตรใหม่ที่ยังไม่มีบัณฑิตรุ่นแรก ปีการศึกษา~2568 จึงเป็นปีที่หลักสูตรวาง ``Achievement Measurement System'' ให้ครบถ้วน เพื่อใช้ในการวัดจริงเมื่อนักศึกษาสำเร็จการศึกษา (คาดปีการศึกษา~2569--2570)

% ------------------------------------------------------------------------------
\criterionsub{4.6 The programme to show that the results of the student assessment are used in a timely manner for providing feedback to help students improve their learning.}

หลักสูตรกำหนดให้อาจารย์ผู้สอนทุกวิชาให้ Feedback ผลการประเมินแก่นักศึกษาอย่างทันเวลา ตามกระบวนการที่กำหนดใน มคอ.3 ดังนี้

\begin{enumerate}
  \item \textbf{Feedback ระหว่างภาคการศึกษา}: อาจารย์คืนผลการทำงาน/แบบทดสอบย่อยภายใน~2 สัปดาห์ หลังการส่งงาน พร้อมคำแนะนำเพื่อปรับปรุง ซึ่งให้เวลานักศึกษาปรับตัวก่อนการสอบกลางภาค
  \item \textbf{Feedback หลังการสอบกลางภาค}: อาจารย์วิเคราะห์ผลรวมของชั้นเรียน แจ้งคะแนนและชี้แจงเฉลย เพื่อให้นักศึกษาเข้าใจจุดอ่อนและวางแผนการเตรียมสอบปลายภาค
  \item \textbf{Feedback สำหรับวิทยานิพนธ์/สารนิพนธ์}: อาจารย์ที่ปรึกษานัดประชุมกับนักศึกษาอย่างน้อย~1 ครั้งต่อเดือน เพื่อติดตามความก้าวหน้าและให้ Feedback อย่างต่อเนื่อง ผ่านระบบอาจารย์ที่ปรึกษา (advisor.uru.ac.th)
\end{enumerate}

กระบวนการ Feedback นี้ช่วยให้นักศึกษาสามารถวางแผนการเรียนรู้และปรับปรุงผลงานได้ก่อนการประเมินรอบถัดไป ลดความเสี่ยงในการไม่บรรลุ CLOs

% ------------------------------------------------------------------------------
\criterionsub{4.7 The programme to show that its student assessment methods are subject to continuous review in order to ensure their validity, reliability, and relevance to the student assessment policy.}

หลักสูตรวางระบบทบทวนและปรับปรุงวิธีการประเมินอย่างต่อเนื่องผ่านกลไกดังนี้

\textbf{กระบวนการทวนสอบรายวิชา}: ตามประกาศแนวปฏิบัติและการกำกับติดตามการส่งผลการประเมิน พ.ศ.~2567 (\hyperlink{ref:AUNQA-4-7-1}{AUNQA-4-7-1}) มหาวิทยาลัยกำหนดให้มีการทวนสอบผลการประเมินของรายวิชา และหลักสูตรมีการจัดประชุมคณะกรรมการบริหารหลักสูตรเพื่อทวนสอบข้อสอบและทวนสอบผลสัมฤทธิ์ปลายภาคเรียน (อ้างอิง~\hyperlink{ref:AUNQA-4-5-1}{AUNQA-4-5-1}) เพื่อตรวจสอบว่าการประเมินสอดคล้องกับ CLOs และเป็นไปตามมาตรฐานที่กำหนด

\textbf{การรายงานผล มคอ.5}: เมื่อสิ้นสุดแต่ละภาคการศึกษา อาจารย์ผู้สอนจัดทำ มคอ.5 ซึ่งประกอบด้วยการวิเคราะห์ผลการประเมิน ปัญหาที่พบ และข้อเสนอปรับปรุง รายงานเหล่านี้เป็นข้อมูลนำเข้าสำหรับการทบทวนวิธีการประเมินในภาคการศึกษาถัดไป

\textbf{การทบทวนโดยบุคคลที่~3 (Third-party Review)}: ระบบทวนสอบผลสัมฤทธิ์ผู้เรียนของมหาวิทยาลัยกำหนดให้มีคณะกรรมการทวนสอบที่ประกอบด้วยบุคลากรที่ไม่ใช่อาจารย์ผู้สอนรายวิชานั้นๆ ดำเนินการตรวจสอบว่าการประเมินในแต่ละรายวิชาสอดคล้องกับ CLOs และ PLOs อย่างแท้จริง กระบวนการนี้ดำเนินการในช่วงสิ้นภาคการศึกษา โดยมีขั้นตอนชัดเจน ดังนี้

\begin{enumerate}
  \item \textbf{ใคร}: คณะกรรมการทวนสอบผลสัมฤทธิ์ผู้เรียน ประกอบด้วย ประธานหลักสูตร (ผศ.ดร.พิศิษฐ์ นาคใจ) และอาจารย์ประจำหลักสูตรที่ไม่ได้สอนวิชานั้น อย่างน้อย~2 คน ทำหน้าที่เป็นบุคคลที่~3
  \item \textbf{เมื่อไหร่}: ดำเนินการทุกสิ้นภาคการศึกษา พร้อมกับการรวบรวมรายงาน มคอ.5 ของทุกรายวิชา
  \item \textbf{พิจารณาอะไร}: (1) ตรวจสอบว่า Rubric และ Marking Scheme วัด CLOs ได้จริง (2) ทบทวนความสอดคล้องระหว่างเกณฑ์ประเมินกับ Learning Level ตาม Bloom's Taxonomy (3) ตรวจสอบว่าผลการประเมินสะท้อน Achievement ได้ตรงตามเจตนา
  \item \textbf{ผลลัพธ์}: บันทึกผลการทวนสอบและข้อเสนอแนะในรายงานประชุม (\hyperlink{ref:AUNQA-4-5-1}{AUNQA-4-5-1}) ซึ่งอาจารย์ผู้สอนนำไปปรับปรุงใน มคอ.3 รอบถัดไป
\end{enumerate}

หลักฐานรอบแรก: การประชุมคณะกรรมการบริหารหลักสูตร ครั้งที่~1/2569 (27~มีนาคม~2569) มีวาระทวนสอบผลสัมฤทธิ์ผู้เรียน ภาคการศึกษาที่~2/2568 ซึ่งเป็นการดำเนินการ Third-party Review รอบแรกของหลักสูตร (อ้างอิง~\hyperlink{ref:AUNQA-4-5-1}{AUNQA-4-5-1})

\textbf{การประสาน Feedback กลับสู่การพัฒนา TLA}: ผลการทบทวนการประเมินเชื่อมโยงกับการปรับปรุงกิจกรรมการเรียนการสอน (R3.6) ผ่านกระบวนการประชุมคณะกรรมการบริหารหลักสูตรประจำภาคการศึกษา ทำให้เกิดวงจร PDCA ระหว่างการสอน (C3) และการประเมิน (C4) อย่างต่อเนื่อง ตัวอย่างที่เกิดขึ้นจริง: วิชา~7015101 รายงาน มคอ.5 ว่านักศึกษา~1 คนมีปัญหาภาษาอังกฤษในการอ่านบทความ $\rightarrow$ แผนปรับปรุง: บูรณาการทักษะภาษาอังกฤษวิชาการเข้าสู่กิจกรรมสัมมนา (7015907) ในภาคถัดไป นี่คือตัวอย่าง PDCA loop ที่ Assessment นำไปสู่การปรับปรุง TLA จริง

\begin{strengthbox}
หลักสูตรมีวิธีประเมินที่หลากหลายและ Constructive Aligned กับ CLOs ทุกวิชา มีระบบอุทธรณ์อย่างเป็นทางการ (\hyperlink{ref:AUNQA-4-2-1}{AUNQA-4-2-1}) ข้อบังคับความก้าวหน้าบัณฑิตศึกษาครบถ้วน Rubric 4 ระดับทุกวิชา PLO Achievement Criteria ครบ 5 PLO ระบบ Feedback 3 ช่วงเวลา และระบบทวนสอบ (บุคคลที่~3) อย่างเป็นทางการ
\end{strengthbox}

\begin{weaknessbox}
PLO Achievement Level ระดับ Graduation จะรายงานได้เมื่อนักศึกษาสำเร็จการศึกษา (คาดปีการศึกษา~2569--2570); บันทึกการทวนสอบ (Third-party Review) รอบแรกยังอยู่ในรูปรายงานประชุมภายใน ยังไม่ได้จัดทำเป็นแบบฟอร์มมาตรฐานรายวิชา
\end{weaknessbox}

\begin{selfassessmentbox}{2}{หลักสูตรมีระบบการประเมินผู้เรียนที่ครบถ้วน: วิธีประเมินหลากหลาย Constructive Aligned กับ CLOs, Rubrics 4 ระดับทุกวิชาบังคับ, ระบบอุทธรณ์เป็นทางการ (\hyperlink{ref:AUNQA-4-2-1}{AUNQA-4-2-1}), Feedback 3 ช่วงเวลา, Third-party Review ดำเนินการรอบแรกแล้ว (27 มี.ค. 2569), CLO Achievement 100\% ทุกวิชาบังคับยืนยันจาก มคอ.5, และ PLO Achievement Criteria ครบ 5 PLO พร้อม Aggregate 100\% ปีการศึกษา 2568 (ระดับรายวิชา; PLO ระดับบัณฑิตรายงานเมื่อสำเร็จการศึกษา คาดปีการศึกษา~2569--2570)}
\end{selfassessmentbox}

\begin{evidencelist}
\evidence{AUNQA-4-1-1}{ตาราง Assessment Methods รายวิชาบังคับ (มคอ.3 ทุกวิชา)\ \href{https://syweb.kidcbc.work/Eviden_aunqa68/\%E0\%B8\%A1\%E0\%B8\%84\%E0\%B8\%AD3/AUN3_4095101.pdf}{\texttt{[PDF]}}}
\evidence{AUNQA-4-2-1}{ประกาศระบบและแนวทางการอุทธรณ์ผลการประเมิน พ.ศ.2567\ \href{https://syweb.kidcbc.work/Eviden_aunqa68/\%E0\%B8\%A3\%E0\%B8\%B0\%E0\%B9\%80\%E0\%B8\%9A\%E0\%B8\%B5\%E0\%B8\%A2\%E0\%B8\%9A\%E0\%B8\%82\%E0\%B9\%89\%E0\%B8\%AD\%E0\%B8\%9A\%E0\%B8\%B1\%E0\%B8\%87\%E0\%B8\%84\%E0\%B8\%B1\%E0\%B8\%9A\%E0\%B9\%81\%E0\%B8\%A5\%E0\%B8\%B0\%E0\%B8\%9B\%E0\%B8\%A3\%E0\%B8\%B0\%E0\%B8\%81\%E0\%B8\%B2\%E0\%B8\%A8\%E0\%B8\%AD\%E0\%B8\%B7\%E0\%B9\%88\%E0\%B8\%99\%E0\%B9\%86/AUNQA-4-2-1_\%E0\%B8\%9B\%E0\%B8\%A3\%E0\%B8\%B0\%E0\%B8\%81\%E0\%B8\%B2\%E0\%B8\%A8\%E0\%B8\%AD\%E0\%B8\%B8\%E0\%B8\%97\%E0\%B8\%98\%E0\%B8\%A3\%E0\%B8\%93\%E0\%B9\%8C\%E0\%B8\%9C\%E0\%B8\%A5\%E0\%B8\%9B\%E0\%B8\%A3\%E0\%B8\%B0\%E0\%B9\%80\%E0\%B8\%A1\%E0\%B8\%B4\%E0\%B8\%99_2567.pdf}{\texttt{[PDF]}}}
\evidence{AUNQA-4-3-1}{ข้อบังคับ มรอ. ว่าด้วยการศึกษาระดับบัณฑิตศึกษา พ.ศ.2566\ \href{https://syweb.kidcbc.work/Eviden_aunqa68/\%E0\%B8\%A3\%E0\%B8\%B0\%E0\%B9\%80\%E0\%B8\%9A\%E0\%B8\%B5\%E0\%B8\%A2\%E0\%B8\%9A\%E0\%B8\%82\%E0\%B9\%89\%E0\%B8\%AD\%E0\%B8\%9A\%E0\%B8\%B1\%E0\%B8\%87\%E0\%B8\%84\%E0\%B8\%B1\%E0\%B8\%9A\%E0\%B9\%81\%E0\%B8\%A5\%E0\%B8\%B0\%E0\%B8\%9B\%E0\%B8\%A3\%E0\%B8\%B0\%E0\%B8\%81\%E0\%B8\%B2\%E0\%B8\%A8\%E0\%B8\%AD\%E0\%B8\%B7\%E0\%B9\%88\%E0\%B8\%99\%E0\%B9\%86/AUNQA-4-3-1_\%E0\%B8\%82\%E0\%B9\%89\%E0\%B8\%AD\%E0\%B8\%9A\%E0\%B8\%B1\%E0\%B8\%87\%E0\%B8\%84\%E0\%B8\%B1\%E0\%B8\%9A\%E0\%B8\%9A\%E0\%B8\%B1\%E0\%B8\%93\%E0\%B8\%91\%E0\%B8\%B4\%E0\%B8\%95\%E0\%B8\%A8\%E0\%B8\%B6\%E0\%B8\%81\%E0\%B8\%A9\%E0\%B8\%B2_2566.pdf}{\texttt{[PDF]}}}
\evidence{AUNQA-4-3-2}{ข้อบังคับ มรอ. ว่าด้วยการศึกษาระดับบัณฑิตศึกษา (ฉบับที่ 2)\ \href{https://syweb.kidcbc.work/Eviden_aunqa68/\%E0\%B8\%A3\%E0\%B8\%B0\%E0\%B9\%80\%E0\%B8\%9A\%E0\%B8\%B5\%E0\%B8\%A2\%E0\%B8\%9A\%E0\%B8\%82\%E0\%B9\%89\%E0\%B8\%AD\%E0\%B8\%9A\%E0\%B8\%B1\%E0\%B8\%87\%E0\%B8\%84\%E0\%B8\%B1\%E0\%B8\%9A\%E0\%B9\%81\%E0\%B8\%A5\%E0\%B8\%B0\%E0\%B8\%9B\%E0\%B8\%A3\%E0\%B8\%B0\%E0\%B8\%81\%E0\%B8\%B2\%E0\%B8\%A8\%E0\%B8\%AD\%E0\%B8\%B7\%E0\%B9\%88\%E0\%B8\%99\%E0\%B9\%86/AUNQA-4-3-2_\%E0\%B8\%82\%E0\%B9\%89\%E0\%B8\%AD\%E0\%B8\%9A\%E0\%B8\%B1\%E0\%B8\%87\%E0\%B8\%84\%E0\%B8\%B1\%E0\%B8\%9A\%E0\%B8\%9A\%E0\%B8\%B1\%E0\%B8\%93\%E0\%B8\%91\%E0\%B8\%B4\%E0\%B8\%95\%E0\%B8\%A8\%E0\%B8\%B6\%E0\%B8\%81\%E0\%B8\%A9\%E0\%B8\%B2\%E0\%B8\%892_2566.pdf}{\texttt{[PDF]}}}
\evidence{AUNQA-4-5-1}{รายงานการประชุมคณะกรรมการบริหารหลักสูตร ครั้งที่ 1/2569 (ทวนสอบผลสัมฤทธิ์ผู้เรียน); รายงานสรุปภาค 2/2568 --- meeting-minutes/report\_meeting\_2569-03-27\_semester2-closeout.pdf\ \href{https://syweb.kidcbc.work/Eviden_aunqa68/\%E0\%B8\%AB\%E0\%B8\%A5\%E0\%B8\%B1\%E0\%B8\%81\%E0\%B8\%90\%E0\%B8\%B2\%E0\%B8\%99\%E0\%B8\%81\%E0\%B8\%B2\%E0\%B8\%A3\%E0\%B8\%9B\%E0\%B8\%A3\%E0\%B8\%B0\%E0\%B8\%8A\%E0\%B8\%B8\%E0\%B8\%A1/01_\%E0\%B8\%A3\%E0\%B8\%B0\%E0\%B8\%94\%E0\%B8\%B1\%E0\%B8\%9A\%E0\%B8\%AB\%E0\%B8\%A5\%E0\%B8\%B1\%E0\%B8\%81\%E0\%B8\%AA\%E0\%B8\%B9\%E0\%B8\%95\%E0\%B8\%A3/\%E0\%B8\%81\%E0\%B8\%A1\%E0\%B8\%A8_2569-03-27_\%E0\%B8\%AA\%E0\%B8\%A3\%E0\%B8\%B8\%E0\%B8\%9B\%E0\%B8\%9C\%E0\%B8\%A5\%E0\%B8\%A0\%E0\%B8\%B2\%E0\%B8\%84\%E0\%B9\%80\%E0\%B8\%A3\%E0\%B8\%B5\%E0\%B8\%A2\%E0\%B8\%992.pdf}{\texttt{[PDF]}}}
\evidence{AUNQA-4-7-1}{ประกาศแนวปฏิบัติและการกำกับติดตามการส่งผลการประเมิน พ.ศ.2567\ \href{https://syweb.kidcbc.work/Eviden_aunqa68/\%E0\%B8\%A3\%E0\%B8\%B0\%E0\%B9\%80\%E0\%B8\%9A\%E0\%B8\%B5\%E0\%B8\%A2\%E0\%B8\%9A\%E0\%B8\%82\%E0\%B9\%89\%E0\%B8\%AD\%E0\%B8\%9A\%E0\%B8\%B1\%E0\%B8\%87\%E0\%B8\%84\%E0\%B8\%B1\%E0\%B8\%9A\%E0\%B9\%81\%E0\%B8\%A5\%E0\%B8\%B0\%E0\%B8\%9B\%E0\%B8\%A3\%E0\%B8\%B0\%E0\%B8\%81\%E0\%B8\%B2\%E0\%B8\%A8\%E0\%B8\%AD\%E0\%B8\%B7\%E0\%B9\%88\%E0\%B8\%99\%E0\%B9\%86/AUNQA-4-7-1_\%E0\%B9\%81\%E0\%B8\%99\%E0\%B8\%A7\%E0\%B8\%9B\%E0\%B8\%8F\%E0\%B8\%B4\%E0\%B8\%9A\%E0\%B8\%B1\%E0\%B8\%95\%E0\%B8\%B4\%E0\%B8\%AA\%E0\%B9\%88\%E0\%B8\%87\%E0\%B8\%9C\%E0\%B8\%A5\%E0\%B8\%9B\%E0\%B8\%A3\%E0\%B8\%B0\%E0\%B9\%80\%E0\%B8\%A1\%E0\%B8\%B4\%E0\%B8\%99_2567.pdf}{\texttt{[PDF]}}}
\evidence{AUNQA-C6-STSAT-2568}{แบบสำรวจความพึงพอใจนักศึกษาและผู้มีส่วนได้ส่วนเสีย CPE\&AI ปีการศึกษา~2568 ($n=4$): คะแนนด้านความชัดเจนเกณฑ์ประเมินและกิจกรรม Practical~5.00/5.00 --- student\_satisfaction\_survey\_2568.csv \href{https://docs.google.com/forms/d/1xpfb6AfYkbOkidJNh6ci95ux4ZLYzjvpUuqARYjhr5k/viewanalytics}{\texttt{[Google Report]}}}
\end{evidencelist}
